\chapter{Solutions to Chapter 13: The Average Causal Effect on the Treated Units and Other Estimands}
\label{sol:chapter13}

\begin{exercisesolution}{Comparing $\tau_\textsc{T}, \tau_\textsc{C},$ and $\tau$}{propensity score 与 CATE 的 covariance}
本题的核心是把 treated population 与 control population 看成对原 population 的重新加权。由 consistency 和 ignorability，
\[
\tau_\textsc{T}
=\mathbb{E}\{Y(1)-Y(0)\mid Z=1\}
=\mathbb{E}\{\tau(X)\mid Z=1\}.
\]
对任意可积函数 $g(X)$，都有
\[
\mathbb{E}\{g(X)\mid Z=1\}
=\frac{\mathbb{E}\{Zg(X)\}}{\mathbb{P}(Z=1)}
=\frac{\mathbb{E}\{e(X)g(X)\}}{e}.
\]
取 $g(X)=\tau(X)$，得到
\[
\tau_\textsc{T}
=\frac{\mathbb{E}\{e(X)\tau(X)\}}{e}.
\]
另一方面，
\[
\tau=\mathbb{E}\{\tau(X)\}.
\]
因此
\begin{align*}
\tau_\textsc{T}-\tau
&=\frac{\mathbb{E}\{e(X)\tau(X)\}}{e}
  -\mathbb{E}\{\tau(X)\} \\
&=\frac{\mathbb{E}\{e(X)\tau(X)\}
  -e\mathbb{E}\{\tau(X)\}}{e} \\
&=\frac{\text{cov}\{e(X),\tau(X)\}}{e}.
\end{align*}

同理，
\[
\mathbb{E}\{g(X)\mid Z=0\}
=\frac{\mathbb{E}[\{1-e(X)\}g(X)]}{1-e}.
\]
于是
\begin{align*}
\tau_\textsc{C}-\tau
&=\frac{\mathbb{E}[\{1-e(X)\}\tau(X)]}{1-e}
  -\mathbb{E}\{\tau(X)\} \\
&=\frac{\mathbb{E}\{\tau(X)\}
  -\mathbb{E}\{e(X)\tau(X)\}
  -(1-e)\mathbb{E}\{\tau(X)\}}{1-e} \\
&=-\frac{\mathbb{E}\{e(X)\tau(X)\}
  -e\mathbb{E}\{\tau(X)\}}{1-e} \\
&=-\frac{\text{cov}\{e(X),\tau(X)\}}{1-e}.
\end{align*}

\paragraph{Remark.}
这个恒等式给出了一个很直接的诊断：如果 treatment 更容易发生在 treatment effect 较大的 covariate 区域，那么 $\text{cov}\{e(X),\tau(X)\}>0$，于是 $\tau_\textsc{T}>\tau>\tau_\textsc{C}$。反过来，如果 treatment 倾向于发生在受益较小的人群中，三个 estimands 的排序会反转。
\end{exercisesolution}

\begin{exercisesolution}{An algebraic fact about a regression estimator for $\tau_\textsc{T}$}{交互 OLS 的 treatment 系数}
本题补全 \cref{eg::att-regression2} 末尾的代数事实。记
\[
\tilde X_i=X_i-\hat{\bar X}(1),\qquad
\hat{\bar X}(1)=n_1^{-1}\sumn Z_iX_i.
\]
考虑所有 units 上的 fully interacted OLS：
\[
Y_i=\gamma_0+\gamma_zZ_i+\gamma_x\tran \tilde X_i
      +\gamma_{zx}\tran Z_i\tilde X_i+\varepsilon_i.
\]
由于模型包含 $Z_i\tilde X_i$，它等价于在 control group 和 treatment group 中分别拟合两条线性回归。对 control group，
\[
Y_i=\gamma_0+\gamma_x\tran \tilde X_i+\varepsilon_i,\qquad Z_i=0;
\]
对 treatment group，
\[
Y_i=(\gamma_0+\gamma_z)+(\gamma_x+\gamma_{zx})\tran \tilde X_i+\varepsilon_i,\qquad Z_i=1.
\]
因此 $\gamma_0$ 是 control regression 在 $\tilde X=0$，也就是 $X=\hat{\bar X}(1)$ 处的 fitted value；而 $\gamma_0+\gamma_z$ 是 treatment regression 在同一点处的 fitted value。

先看 treatment group。因为 treatment group 中
\[
n_1^{-1}\sumn Z_i\tilde X_i=0,
\]
含 intercept 的 OLS 正规方程给出
\[
\hat\gamma_0+\hat\gamma_z
=n_1^{-1}\sumn Z_iY_i
=\hat{\bar Y}(1).
\]
再看 control group。只在 control units 上回归 $Y$ 对 intercept 与 $X$，得到 \cref{eg::att-regression2} 中的
\[
\hat\beta_{0\mid 0},\qquad \hat\beta_{x\mid 0}.
\]
同一条 control fitted line 用中心化 covariate 表示时，在 $\tilde X=0$ 处的值为
\[
\hat\gamma_0
=\hat\beta_{0\mid 0}+\hat\beta_{x\mid 0}\tran \hat{\bar X}(1).
\]
于是 fully interacted OLS 中 $Z$ 的系数为
\begin{align*}
\hat\gamma_z
&=(\hat\gamma_0+\hat\gamma_z)-\hat\gamma_0 \\
&=\hat{\bar Y}(1)
  -\hat\beta_{0\mid 0}
  -\hat\beta_{x\mid 0}\tran \hat{\bar X}(1).
\end{align*}
这正是 \cref{eg::att-regression2} 中的 estimator
\[
\hat{\tau}_\textsc{T}
=\hat{\bar Y}(1)
  -\hat\beta_{0\mid 0}
  -\hat\beta_{x\mid 0}\tran \hat{\bar X}(1).
\]

\paragraph{Remark.}
中心化点必须是 treated covariate mean。这样 $Z$ 的系数才是在 treated population 的 covariate profile 上比较 treated regression surface 与 control regression surface，而不是在 overall mean 或 control mean 上比较。
\end{exercisesolution}

\begin{exercisesolution}{Simulation studies for $\tau_\textsc{T}$}{DR estimator 的有限样本表现}
本题要求仿照 \cref{chapter::doubly-robust} 对 $\tau$ 的 simulation studies，为 $\tau_\textsc{T}$ 比较 outcome regression、HT、Hajek 和 doubly robust estimators。下面给出一个可复现实验。生成机制为
\[
X\sim N(0,1),\qquad e(X)=\frac{\exp(-0.2+0.8X)}{1+\exp(-0.2+0.8X)},
\]
\[
Z\mid X\sim \operatorname{Bernoulli}\{e(X)\},
\qquad
Y(0)=1+X+\frac{1}{2}X^2+\varepsilon,\quad \varepsilon\sim N(0,1),
\]
\[
Y(1)=Y(0)+1+\frac{1}{2}X.
\]
因此
\[
\tau(X)=1+\frac{1}{2}X,\qquad
\tau_\textsc{T}=\frac{\mathbb{E}\{e(X)\tau(X)\}}{\mathbb{E}\{e(X)\}}.
\]
用 $10^6$ 次 Monte Carlo integration 近似得到
\[
\mathbb{E}\{e(X)\}\approx 0.456,\qquad
\tau_\textsc{T}\approx 1.191.
\]

每次 simulation 取 $n=1000$，重复 $300$ 次；每次用 $80$ 个 bootstrap samples 估计 standard error。propensity score model 正确时拟合 logistic regression $Z\sim X$，错误时只拟合 intercept。outcome model 正确时在 control units 中拟合 $Y\sim X+X^2$，错误时只拟合 $Y\sim X$。得到如下结果：
\[
\begin{array}{cccrrrr}
\hline
\text{PS} & \text{OM} & \text{Estimator} & \text{Mean} & \text{Bias} & \text{SD} & \text{Mean boot SE}\\
\hline
\text{correct} & \text{correct} & \text{reg}   & 1.187 & -0.004 & 0.079 & 0.077\\
\text{correct} & \text{correct} & \text{HT}    & 1.177 & -0.014 & 0.177 & 0.163\\
\text{correct} & \text{correct} & \text{Hajek} & 1.180 & -0.010 & 0.152 & 0.139\\
\text{correct} & \text{correct} & \text{DR}    & 1.186 & -0.004 & 0.079 & 0.077\\
\text{wrong} & \text{correct} & \text{reg}   & 1.192 &  0.001 & 0.077 & 0.077\\
\text{wrong} & \text{correct} & \text{HT}    & 1.915 &  0.724 & 0.120 & 0.115\\
\text{wrong} & \text{correct} & \text{Hajek} & 1.915 &  0.724 & 0.120 & 0.115\\
\text{wrong} & \text{correct} & \text{DR}    & 1.192 &  0.001 & 0.077 & 0.077\\
\text{correct} & \text{wrong} & \text{reg}   & 1.444 &  0.253 & 0.104 & 0.106\\
\text{correct} & \text{wrong} & \text{HT}    & 1.198 &  0.008 & 0.155 & 0.158\\
\text{correct} & \text{wrong} & \text{Hajek} & 1.199 &  0.009 & 0.134 & 0.135\\
\text{correct} & \text{wrong} & \text{DR}    & 1.199 &  0.008 & 0.116 & 0.117\\
\text{wrong} & \text{wrong} & \text{reg}   & 1.439 &  0.249 & 0.105 & 0.107\\
\text{wrong} & \text{wrong} & \text{HT}    & 1.912 &  0.722 & 0.110 & 0.114\\
\text{wrong} & \text{wrong} & \text{Hajek} & 1.912 &  0.722 & 0.110 & 0.114\\
\text{wrong} & \text{wrong} & \text{DR}    & 1.439 &  0.249 & 0.105 & 0.107\\
\hline
\end{array}
\]

结果与 \cref{thm::dr-att} 的 double robustness 完全一致。当 outcome model 正确时，regression estimator 与 DR estimator 几乎无偏；即使 propensity score model 错误，DR estimator 仍然保持无偏。当 propensity score model 正确而 outcome model 错误时，HT、Hajek 和 DR estimators 都接近真实 $\tau_\textsc{T}$；其中 DR 的 variability 小于纯 weighting estimators。当两个 nuisance models 都错时，DR estimator 不再有保护，它退化到错误 outcome regression 的偏差水平。

复现代码如下。
\begin{lstlisting}
set.seed(20260606)

expit <- function(x) 1 / (1 + exp(-x))

true_att_mc <- function(N = 1000000) {
  x <- rnorm(N)
  e <- expit(-0.2 + 0.8 * x)
  tau <- 1 + 0.5 * x
  c(e = mean(e), ATT = mean(e * tau) / mean(e))
}

att_est <- function(z, y, x, ps = c("correct", "wrong"),
                    om = c("correct", "wrong")) {
  ps <- match.arg(ps)
  om <- match.arg(om)
  n <- length(z)
  n1 <- sum(z)
  dat <- data.frame(z = z, y = y, x = x)

  if (ps == "correct") {
    phat <- glm(z ~ x, data = dat, family = binomial())$fitted.values
  } else {
    phat <- glm(z ~ 1, data = dat, family = binomial())$fitted.values
  }
  phat <- pmin(0.995, pmax(0.005, phat))
  odds <- phat / (1 - phat)

  if (om == "correct") {
    m0 <- glm(y ~ x + I(x^2), data = dat, weights = 1 - z,
              family = gaussian())$fitted.values
  } else {
    m0 <- glm(y ~ x, data = dat, weights = 1 - z,
              family = gaussian())$fitted.values
  }

  reg <- mean(y[z == 1]) - mean(m0[z == 1])
  ht <- mean(y[z == 1]) - mean(odds * (1 - z) * y) * n / n1
  hajek <- mean(y[z == 1]) -
           sum(odds * (1 - z) * y) / sum(odds * (1 - z))
  dr <- reg - mean(odds * (1 - z) * (y - m0)) * n / n1
  c(reg = reg, HT = ht, Hajek = hajek, DR = dr)
}

one_run <- function(n = 1000, ps = "correct",
                    om = "correct", B = 80) {
  x <- rnorm(n)
  e <- expit(-0.2 + 0.8 * x)
  z <- rbinom(n, 1, e)
  y0 <- 1 + x + 0.5 * x^2 + rnorm(n)
  y1 <- y0 + 1 + 0.5 * x
  y <- z * y1 + (1 - z) * y0

  est <- att_est(z, y, x, ps, om)
  boot <- replicate(B, {
    id <- sample.int(n, n, replace = TRUE)
    att_est(z[id], y[id], x[id], ps, om)
  })
  c(est, bootse = apply(boot, 1, sd))
}

truth <- true_att_mc()
settings <- expand.grid(ps = c("correct", "wrong"),
                        om = c("correct", "wrong"),
                        stringsAsFactors = FALSE)
nsim <- 300
B <- 80

res <- lapply(seq_len(nrow(settings)), function(j) {
  ps <- settings$ps[j]
  om <- settings$om[j]
  mat <- replicate(nsim, one_run(ps = ps, om = om, B = B))
  est <- t(mat[1:4, ])
  bse <- t(mat[5:8, ])
  data.frame(ps = ps, om = om, estimator = colnames(est),
             mean = colMeans(est),
             bias = colMeans(est) - truth["ATT"],
             sd = apply(est, 2, sd),
             mean_boot_se = colMeans(bse))
})

out <- do.call(rbind, res)
round(out, 3)
\end{lstlisting}

\paragraph{Remark.}
这个 simulation 也说明 ATT 问题中 HT estimator 的不稳定性。即使 propensity score model 正确，HT 的 Monte Carlo standard deviation 仍明显大于 regression 和 DR estimators；Hajek normalization 降低了一部分波动，但不能替代一个有预测力的 outcome model。
\end{exercisesolution}

\begin{exercisesolution}{An alternative form of the doubly robust estimator for $\tau_\textsc{T}$}{Robins 型 normalized residual correction}
记
\[
o_\alpha(X)=\frac{e(X,\alpha)}{1-e(X,\alpha)},\qquad
m_0(X)=\mu_0(X,\beta_0).
\]
题目中的 estimator 可写成
\[
\tilde{\mu}_{0\textsc{T}}^\textup{dr2}
=\frac{\mathbb{E}[o_\alpha(X)(1-Z)\{Y-m_0(X)\}]}
        {\mathbb{E}\{o_\alpha(X)(1-Z)\}}
 +\frac{\mathbb{E}\{Zm_0(X)\}}{e}.
\]

先假设 propensity score model 正确，即 $e(X,\alpha)=e(X)$。分母为
\[
\mathbb{E}\{o(X)(1-Z)\}
=\mathbb{E}\left[
\frac{e(X)}{1-e(X)}\mathbb{E}(1-Z\mid X)
\right]
=\mathbb{E}\{e(X)\}=e.
\]
分子为
\begin{align*}
\mathbb{E}[o(X)(1-Z)\{Y-m_0(X)\}]
&=\mathbb{E}\left[
\frac{e(X)}{1-e(X)}
\mathbb{E}\{(1-Z)(Y(0)-m_0(X))\mid X\}
\right] \\
&=\mathbb{E}\left[e(X)\{\mu_0(X)-m_0(X)\}\right].
\end{align*}
因此
\begin{align*}
\tilde{\mu}_{0\textsc{T}}^\textup{dr2}
&=\frac{\mathbb{E}[e(X)\{\mu_0(X)-m_0(X)\}]}{e}
  +\frac{\mathbb{E}\{e(X)m_0(X)\}}{e} \\
&=\frac{\mathbb{E}\{e(X)\mu_0(X)\}}{e}
=\mathbb{E}\{Y(0)\mid Z=1\}.
\end{align*}

再假设 outcome model 正确，即 $m_0(X)=\mu_0(X)$，但不要求 propensity score model 正确。条件于 $X$，
\[
\mathbb{E}\{(1-Z)(Y-m_0(X))\mid X\}
=\mathbb{E}\{(1-Z)(Y(0)-\mu_0(X))\mid X\}=0.
\]
所以 normalized residual correction 的分子为 $0$。同时
\[
\frac{\mathbb{E}\{Z\mu_0(X)\}}{e}
=\frac{\mathbb{E}\{e(X)\mu_0(X)\}}{e}
=\mathbb{E}\{Y(0)\mid Z=1\}.
\]
这证明了 $\tilde{\mu}_{0\textsc{T}}^\textup{dr2}$ 的 double robustness。

样本 analog 为
\[
\hat{\mu}_{0\textsc{T}}^\textup{dr2}
=
\frac{\sumn \hat o_i(1-Z_i)\{Y_i-\hat m_{0i}\}}
     {\sumn \hat o_i(1-Z_i)}
+\frac{1}{n_1}\sumn Z_i\hat m_{0i},
\]
其中
\[
\hat o_i=\frac{e(X_i,\hat\alpha)}{1-e(X_i,\hat\alpha)},\qquad
\hat m_{0i}=\mu_0(X_i,\hat\beta_0).
\]
于是
\[
\hat\tau_\textsc{T}^\textup{dr2}
=\hat{\bar Y}(1)-\hat{\mu}_{0\textsc{T}}^\textup{dr2}.
\]

\paragraph{Remark.}
这个形式与 \cref{def::dr-att} 的区别在于 residual correction 被 odds weights 的总量 normalize。它类似于从 HT 到 Hajek 的改动：population 层面不改变 double robustness，finite sample 中常常减轻权重总量偏离目标带来的不稳定性。
\end{exercisesolution}

\begin{exercisesolution}{Average causal effect on the control units}{ATC 的识别、IPW 与 DR 形式}
对 $\tau_\textsc{C}$，需要识别的 counterfactual mean 是 $\mathbb{E}\{Y(1)\mid Z=0\}$。与 \cref{assume::ignorability-y0} 对称，假设
\[
Z\ind Y(1)\mid X,\qquad e(X)>0.
\]
则
\begin{align*}
\mathbb{E}\{Y(1)\mid Z=0\}
&=\mathbb{E}\left[\mathbb{E}\{Y(1)\mid Z=0,X\}\mid Z=0\right] \\
&=\mathbb{E}\left[\mathbb{E}\{Y(1)\mid Z=1,X\}\mid Z=0\right] \\
&=\mathbb{E}\left\{\mathbb{E}(Y\mid Z=1,X)\mid Z=0\right\}.
\end{align*}
因此 outcome-regression identification formula 为
\[
\tau_\textsc{C}
=\mathbb{E}\left\{\mathbb{E}(Y\mid Z=1,X)\mid Z=0\right\}
 -\mathbb{E}(Y\mid Z=0).
\]
这正是 \cref{eq::att-identification} 的 control-population analog。

IPW 形式也完全对称。令 $1-e=\mathbb{P}(Z=0)$，则
\begin{align*}
\mathbb{E}\left\{
\frac{1-e(X)}{1-e}\frac{Z}{e(X)}Y
\right\}
&=\mathbb{E}\left[
\frac{1-e(X)}{(1-e)e(X)}
\mathbb{E}\{ZY(1)\mid X\}
\right] \\
&=\mathbb{E}\left[
\frac{1-e(X)}{(1-e)e(X)}
e(X)\mu_1(X)
\right] \\
&=\frac{\mathbb{E}[\{1-e(X)\}\mu_1(X)]}{1-e}
=\mathbb{E}\{Y(1)\mid Z=0\}.
\end{align*}
于是
\[
\tau_\textsc{C}
=\mathbb{E}\left\{
\frac{1-e(X)}{1-e}\frac{Z}{e(X)}Y
\right\}
-\mathbb{E}(Y\mid Z=0).
\]
若记 inverse odds
\[
r(X)=\frac{1-e(X)}{e(X)},
\]
则 HT 和 Hajek estimators 可写为
\[
\hat\tau_\textsc{C}^\textup{ht}
=n_0^{-1}\sumn \hat r(X_i)Z_iY_i-\hat{\bar Y}(0),
\]
\[
\hat\tau_\textsc{C}^\textup{hajek}
=\frac{\sumn \hat r(X_i)Z_iY_i}{\sumn \hat r(X_i)Z_i}
 -\hat{\bar Y}(0).
\]

最后给出 DR estimator。设 $\mu_1(X,\beta_1)$ 是 treatment outcome model，$r(X,\alpha)=\{1-e(X,\alpha)\}/e(X,\alpha)$。population version 为
\[
\tilde\mu_{1\textsc{C}}^\textup{dr}
=\mathbb{E}\left[
r(X,\alpha)Z\{Y-\mu_1(X,\beta_1)\}
+(1-Z)\mu_1(X,\beta_1)
\right]/(1-e),
\]
并定义
\[
\tilde\tau_\textsc{C}^\textup{dr}
=\tilde\mu_{1\textsc{C}}^\textup{dr}-\mathbb{E}(Y\mid Z=0).
\]
样本 analog 为
\[
\hat\mu_{1\textsc{C}}^\textup{dr}
=\frac{1}{n_0}\sumn
\left[
\hat r_iZ_i\{Y_i-\hat m_{1i}\}
+(1-Z_i)\hat m_{1i}
\right],
\]
\[
\hat\tau_\textsc{C}^\textup{dr}
=\hat\mu_{1\textsc{C}}^\textup{dr}-\hat{\bar Y}(0).
\]

\paragraph{Remark.}
ATT 和 ATC 的公式不是简单地把符号 $1$ 与 $0$ 互换；更重要的是目标 population 变了。ATT 使用 treated covariate distribution，ATC 使用 control covariate distribution。因此两者在 treatment effect heterogeneity 下回答的是不同实际问题。
\end{exercisesolution}

\begin{exercisesolution}{Estimating individual effect and conditional average causal effect}{IPW pseudo-outcome 的无偏性}
先考虑 randomized assignment。由 consistency，
\[
Y_i=Z_iY_i(1)+(1-Z_i)Y_i(0).
\]
题目中的 pseudo-outcome 为
\[
\delta_i=\frac{Z_iY_i}{e}-\frac{(1-Z_i)Y_i}{1-e}
=\frac{Z_iY_i(1)}{e}-\frac{(1-Z_i)Y_i(0)}{1-e}.
\]
因为 $Z_i\ind\{Y_i(1),Y_i(0)\}$，条件于 potential outcomes 有
\begin{align*}
\mathbb{E}\{\delta_i\mid Y_i(1),Y_i(0)\}
&=\frac{\mathbb{E}(Z_i)}{e}Y_i(1)
  -\frac{\mathbb{E}(1-Z_i)}{1-e}Y_i(0)\\
&=Y_i(1)-Y_i(0)
=\tau_i.
\end{align*}
所以
\[
\mathbb{E}(\delta_i-\tau_i)=0.
\]
再取 unconditional expectation，
\[
\mathbb{E}(\delta_i)=\mathbb{E}(\tau_i)
=\mathbb{E}\{Y(1)-Y(0)\}
=\tau.
\]

再考虑 observational study 中的 ignorability。此时
\[
\delta_i
=\frac{Z_iY_i(1)}{e(X_i)}
 -\frac{(1-Z_i)Y_i(0)}{1-e(X_i)}.
\]
由 $Z_i\ind\{Y_i(1),Y_i(0)\}\mid X_i$，条件于 $X_i,Y_i(1),Y_i(0)$，
\begin{align*}
\mathbb{E}\{\delta_i\mid X_i,Y_i(1),Y_i(0)\}
&=\frac{\mathbb{E}(Z_i\mid X_i)}{e(X_i)}Y_i(1)
 -\frac{\mathbb{E}(1-Z_i\mid X_i)}{1-e(X_i)}Y_i(0)\\
&=Y_i(1)-Y_i(0)
=\tau_i.
\end{align*}
因此
\[
\mathbb{E}(\delta_i-\tau_i)=0.
\]
同时，条件于 $X_i$，
\begin{align*}
\mathbb{E}(\delta_i\mid X_i)
&=\mathbb{E}\{Y_i(1)-Y_i(0)\mid X_i\}\\
&=\tau(X_i).
\end{align*}
所以
\[
\mathbb{E}\{\delta_i-\tau(X_i)\}=0.
\]
进一步，
\[
\mathbb{E}(\delta_i)
=\mathbb{E}\{\mathbb{E}(\delta_i\mid X_i)\}
=\mathbb{E}\{\tau(X_i)\}
=\tau.
\]

\paragraph{Remark.}
$\delta_i$ 是 unbiased predictor，不是 reliable individual-effect estimator。单个 realized unit 仍然只观测一个 potential outcome，另一个 potential outcome 被 inverse-probability term 以极高方差的方式补偿。因此这种 pseudo-outcome 更适合作为平均或条件平均效应学习中的构造工具，而不是直接解释某个个体的真实 causal effect。
\end{exercisesolution}

\begin{exercisesolution}{General estimand and $(\tau_\textsc{T}, \tau_\textsc{C})$}{一般加权 estimand 的两个特例}
一般 estimand 定义为
\[
\tau^h=\frac{\mathbb{E}\{h(X)\tau(X)\}}{\mathbb{E}\{h(X)\}},
\qquad
\tau(X)=\mathbb{E}\{Y(1)-Y(0)\mid X\}.
\]
当 $h(X)=e(X)$ 时，
\[
\mathbb{E}\{h(X)\}=\mathbb{E}\{e(X)\}=e.
\]
并且
\[
\mathbb{E}\{e(X)\tau(X)\}
=\mathbb{E}\{Z\tau(X)\}.
\]
因此
\[
\tau^h
=\frac{\mathbb{E}\{Z\tau(X)\}}{e}
=\mathbb{E}\{\tau(X)\mid Z=1\}.
\]
在 unconfoundedness 下，
\[
\mathbb{E}\{\tau(X)\mid Z=1\}
=\mathbb{E}\{Y(1)-Y(0)\mid Z=1\}
=\tau_\textsc{T}.
\]

当 $h(X)=1-e(X)$ 时，
\[
\tau^h
=\frac{\mathbb{E}[\{1-e(X)\}\tau(X)]}{1-e}
=\mathbb{E}\{\tau(X)\mid Z=0\}
=\tau_\textsc{C}.
\]

\paragraph{Remark.}
\cref{sec::other-estmand-overlap} 的统一公式把 target population 的选择显式化了。选择 $h(X)=1$ 是 combined population；选择 $h(X)=e(X)$ 是 treated population；选择 $h(X)=1-e(X)$ 是 control population。
\end{exercisesolution}

\begin{exercisesolution}{More on $\tau_\textsc{O}$}{overlap population 的 treated/control 表示}
overlap estimand 定义为
\[
\tau_\textsc{O}
=\frac{\mathbb{E}[e(X)\{1-e(X)\}\tau(X)]}
       {\mathbb{E}[e(X)\{1-e(X)\}]}.
\]
先看 treated-population 表示。对 numerator，
\begin{align*}
\mathbb{E}[\{1-e(X)\}\tau(X)\mid Z=1]
&=\frac{\mathbb{E}[Z\{1-e(X)\}\tau(X)]}{e}\\
&=\frac{\mathbb{E}[e(X)\{1-e(X)\}\tau(X)]}{e}.
\end{align*}
对 denominator，
\[
\mathbb{E}\{1-e(X)\mid Z=1\}
=\frac{\mathbb{E}[e(X)\{1-e(X)\}]}{e}.
\]
两者相除，公共因子 $e^{-1}$ 抵消，所以
\[
\tau_\textsc{O}
=\frac{\mathbb{E}[\{1-e(X)\}\tau(X)\mid Z=1]}
       {\mathbb{E}\{1-e(X)\mid Z=1\}}.
\]

control-population 表示同理。因为
\[
\mathbb{E}\{e(X)\tau(X)\mid Z=0\}
=\frac{\mathbb{E}[\{1-e(X)\}e(X)\tau(X)]}{1-e},
\]
且
\[
\mathbb{E}\{e(X)\mid Z=0\}
=\frac{\mathbb{E}[\{1-e(X)\}e(X)]}{1-e},
\]
所以
\[
\tau_\textsc{O}
=\frac{\mathbb{E}\{e(X)\tau(X)\mid Z=0\}}
       {\mathbb{E}\{e(X)\mid Z=0\}}.
\]

\paragraph{Remark.}
这两个表示说明 overlap estimand 可以从 treated side 或 control side 看。站在 treated side，它下调 $e(X)$ 接近 $1$ 的 treated units；站在 control side，它下调 $e(X)$ 接近 $0$ 的 control units。两种说法都在强调同一件事：overlap population 聚焦于 treatment assignment 最不确定的区域。
\end{exercisesolution}

\begin{exercisesolution}{IPW for the general estimand}{\cref{thm::weighted-estimand-li} 的证明}
在 ignorability 和 overlap 下，
\[
Z\ind \{Y(1),Y(0)\}\mid X,\qquad 0<e(X)<1.
\]
由 consistency，
\[
ZY=ZY(1),\qquad (1-Z)Y=(1-Z)Y(0).
\]
因此
\begin{align*}
\mathbb{E}\left\{\frac{ZYh(X)}{e(X)}\right\}
&=\mathbb{E}\left[
\mathbb{E}\left\{\frac{ZY(1)h(X)}{e(X)}\mid X\right\}
\right] \\
&=\mathbb{E}\left[
\frac{h(X)}{e(X)}
\mathbb{E}(Z\mid X)
\mathbb{E}\{Y(1)\mid X\}
\right] \\
&=\mathbb{E}\{h(X)\mu_1(X)\}.
\end{align*}
同理，
\[
\mathbb{E}\left\{\frac{(1-Z)Yh(X)}{1-e(X)}\right\}
=\mathbb{E}\{h(X)\mu_0(X)\}.
\]
两式相减得
\[
\mathbb{E}\left\{
\frac{ZYh(X)}{e(X)}
-\frac{(1-Z)Yh(X)}{1-e(X)}
\right\}
=\mathbb{E}\{h(X)(\mu_1(X)-\mu_0(X))\}.
\]
而
\[
\mu_1(X)-\mu_0(X)
=\mathbb{E}\{Y(1)-Y(0)\mid X\}
=\tau(X).
\]
于是
\[
\frac{
\mathbb{E}\left\{
ZYh(X)/e(X)-(1-Z)Yh(X)/\{1-e(X)\}
\right\}
}{
\mathbb{E}\{h(X)\}
}
=\frac{\mathbb{E}\{h(X)\tau(X)\}}{\mathbb{E}\{h(X)\}}
=\tau^h,
\]
即 \cref{thm::weighted-estimand-li} 成立。

\paragraph{Remark.}
这个证明把 estimand weighting 和 treatment-assignment weighting 分开了。$h(X)$ 决定我们想平均到哪个 population；$1/e(X)$ 与 $1/\{1-e(X)\}$ 则只是把 observed treated/control outcomes 还原到这个目标 population 的工具。
\end{exercisesolution}

\begin{exercisesolution}{Doubly robust estimation for general estimand}{固定 $h(X)$ 时的乘积偏差}
令
\[
m_1(X)=\mu_1(X,\beta_1),\qquad
m_0(X)=\mu_0(X,\beta_0),\qquad
e_\alpha(X)=e(X,\alpha).
\]
先证明 treatment mean 部分。由 consistency 和 ignorability，
\begin{align*}
\tilde{\mu}_1^{h,\textup{dr}}
&=\mathbb{E}\left[
\frac{Zh(X)\{Y(1)-m_1(X)\}}{e_\alpha(X)}
+h(X)m_1(X)
\right] \\
&=\mathbb{E}\left[
\frac{h(X)e(X)}{e_\alpha(X)}
\{\mu_1(X)-m_1(X)\}
+h(X)m_1(X)
\right].
\end{align*}
所以
\begin{align*}
\tilde{\mu}_1^{h,\textup{dr}}-\mathbb{E}\{h(X)Y(1)\}
&=\mathbb{E}\left[
h(X)\left\{\frac{e(X)}{e_\alpha(X)}-1\right\}
\{\mu_1(X)-m_1(X)\}
\right] \\
&=\mathbb{E}\left[
h(X)\frac{e(X)-e_\alpha(X)}{e_\alpha(X)}
\{\mu_1(X)-m_1(X)\}
\right].
\end{align*}
因此，如果 $e_\alpha(X)=e(X)$ 或 $m_1(X)=\mu_1(X)$，上式为 $0$，即
\[
\tilde{\mu}_1^{h,\textup{dr}}=\mathbb{E}\{h(X)Y(1)\}.
\]

control mean 部分完全平行：
\begin{align*}
\tilde{\mu}_0^{h,\textup{dr}}
&=\mathbb{E}\left[
\frac{(1-Z)h(X)\{Y(0)-m_0(X)\}}{1-e_\alpha(X)}
+h(X)m_0(X)
\right] \\
&=\mathbb{E}\left[
\frac{h(X)\{1-e(X)\}}{1-e_\alpha(X)}
\{\mu_0(X)-m_0(X)\}
+h(X)m_0(X)
\right].
\end{align*}
于是
\begin{align*}
\tilde{\mu}_0^{h,\textup{dr}}-\mathbb{E}\{h(X)Y(0)\}
&=\mathbb{E}\left[
h(X)\left\{\frac{1-e(X)}{1-e_\alpha(X)}-1\right\}
\{\mu_0(X)-m_0(X)\}
\right] \\
&=\mathbb{E}\left[
h(X)\frac{e_\alpha(X)-e(X)}{1-e_\alpha(X)}
\{\mu_0(X)-m_0(X)\}
\right].
\end{align*}
如果 $e_\alpha(X)=e(X)$ 或 $m_0(X)=\mu_0(X)$，则
\[
\tilde{\mu}_0^{h,\textup{dr}}=\mathbb{E}\{h(X)Y(0)\}.
\]

最后，若 propensity score model 正确，则上面两个 weighted means 都正确；若 propensity score model 可能错误但两个 outcome models 都正确，则上面两个 weighted means 也都正确。因此在题目给出的两类条件之一成立时，
\[
\frac{\tilde{\mu}_1^{h,\textup{dr}}-\tilde{\mu}_0^{h,\textup{dr}}}
     {\mathbb{E}\{h(X)\}}
=\frac{\mathbb{E}[h(X)\{Y(1)-Y(0)\}]}
       {\mathbb{E}\{h(X)\}}
=\tau^h.
\]

\paragraph{Remark.}
这里的 double robustness 是针对固定的 $h(X)$ 而言的。若 $h(X)$ 本身包含未知 propensity score，例如 ATT 的 $h(X)=e(X)$、ATC 的 $h(X)=1-e(X)$ 或 overlap estimand 的 $h(X)=e(X)\{1-e(X)\}$，那么估计 $h$ 的误差会进入 target population 的定义，不能直接套用本题的固定-$h$ 结论。
\end{exercisesolution}

\begin{exercisesolution}{Analyzing a dataset from the Karolinska Institute}{二元结局下的 ATT 分析框架}
本题要求重新考察 \cref{hw::Karolinska-dr}，但目标 estimand 改为 $\tau_\textsc{T}$。当前本地工程中没有找到 \ri{karolinska.txt}，因此这里不编造 point estimates。下面给出完整可复现代码；将数据文件放在运行目录后，它会输出 \cref{sec::dr-att-example} 中五个 ATT estimators 及 bootstrap standard errors。

\begin{lstlisting}
ATT_est <- function(z, y, x, out.family = binomial(),
                    Utruncps = 1) {
  n <- length(z)
  n1 <- sum(z)
  x <- as.matrix(x)

  pscore <- glm(z ~ ., data = data.frame(z = z, x),
                family = binomial())$fitted.values
  pscore <- pmin(Utruncps, pscore)
  pscore <- pmin(1 - 1e-6, pmax(1e-6, pscore))
  odds <- pscore / (1 - pscore)

  outcome0 <- glm(y ~ ., data = data.frame(y = y, x),
                  weights = 1 - z,
                  family = out.family)$fitted.values

  reg0 <- lm(y ~ z + x)$coef["z"]
  reg <- mean(y[z == 1]) - mean(outcome0[z == 1])
  HT <- mean(y[z == 1]) - mean(odds * (1 - z) * y) * n / n1
  Hajek <- mean(y[z == 1]) -
           sum(odds * (1 - z) * y) / sum(odds * (1 - z))
  DR <- reg - mean(odds * (1 - z) * (y - outcome0)) * n / n1

  c(reg0 = reg0, reg = reg, HT = HT, Hajek = Hajek, DR = DR)
}

OS_ATT <- function(z, y, x, n.boot = 2000,
                   out.family = binomial(),
                   Utruncps = 1) {
  point.est <- ATT_est(z, y, x, out.family, Utruncps)
  n <- length(z)
  x <- as.matrix(x)

  boot.est <- replicate(n.boot, {
    id <- sample.int(n, n, replace = TRUE)
    ATT_est(z[id], y[id], x[id, , drop = FALSE],
            out.family, Utruncps)
  })

  boot.se <- apply(boot.est, 1, sd)
  out <- rbind(est = point.est, se = boot.se)
  round(out, 3)
}

balance_table <- function(z, x) {
  x <- as.data.frame(x)
  ans <- lapply(names(x), function(v) {
    xv <- x[[v]]
    c(mean_treated = mean(xv[z == 1]),
      mean_control = mean(xv[z == 0]),
      diff = mean(xv[z == 1]) - mean(xv[z == 0]))
  })
  out <- do.call(rbind, ans)
  rownames(out) <- names(x)
  round(out, 3)
}

karolinska <- read.table("karolinska.txt", header = TRUE)
z <- karolinska$hvdiag
y <- 1 - (karolinska$year.survival == 1)
x <- as.matrix(karolinska[, c(3, 4, 5)])

table(z)
table(y, z)
balance_table(z, x)

ps <- glm(z ~ x, family = binomial())$fitted.values
round(c(min = min(ps),
        q05 = quantile(ps, 0.05),
        median = median(ps),
        q95 = quantile(ps, 0.95),
        max = max(ps)), 3)

OS_ATT(z, y, x, n.boot = 2000,
       out.family = binomial(), Utruncps = 1)
OS_ATT(z, y, x, n.boot = 2000,
       out.family = binomial(), Utruncps = 0.95)
OS_ATT(z, y, x, n.boot = 2000,
       out.family = binomial(), Utruncps = 0.90)
\end{lstlisting}

解释结果时应当把 $\tau_\textsc{T}$ 与 \cref{hw::Karolinska-dr} 中的 $\tau$ 区分开。$\tau$ 估计的是 combined population 上 large volume hospital diagnosis 的平均效应；$\tau_\textsc{T}$ 估计的是实际在 large volume hospital 诊断的患者，如果他们改在 small volume hospital 诊断，风险会如何变化。若 treated patients 的 age、rural status 或 gender 分布不同于 overall population，那么 $\tau_\textsc{T}$ 与 $\tau$ 可以有实质差异。

报告时至少应包含四项诊断。第一，报告 crude risk difference，但只把它作为 descriptive contrast。第二，报告 estimated propensity scores 的分布，尤其是 $e(X)$ 接近 $1$ 的 treated units；ATT 的 control weights 为 $e(X)/\{1-e(X)\}$，这一区域会放大 control outcomes。第三，比较 HT 与 Hajek；若二者差异很大，说明 odds weights 的总量不稳定。第四，比较 regression、Hajek 和 DR；若 DR 接近 regression 而远离 weighting，说明 outcome model 在主导结论；若 DR 接近 weighting 而远离 regression，说明 control outcome extrapolation 可能较弱。

\paragraph{Remark.}
在 ATT 分析中，overlap 的关注点是不一样的。我们不要求每个 control-like unit 都有 treated counterpart；我们要求 treated covariate profile 能在 control group 中找到足够相近的对照。对 Karolinska 这样样本量只有 $158$ 的数据，truncation sensitivity 和 bootstrap instability 本身就是需要报告的实证发现。
\end{exercisesolution}
